\documentclass[11pt,a4paper]{article}
\usepackage[utf8]{inputenc}
\usepackage[T1]{fontenc}
\usepackage{lmodern}
\usepackage{amsmath,amssymb,amsfonts,amsthm}
\usepackage{geometry}
\usepackage{booktabs}
\usepackage{cite}
\usepackage{microtype}
\usepackage{hyperref}

% Layout: Saubere Randeinstellung, kein Overfull
\geometry{
    a4paper,
    top=2.2cm,
    bottom=2.3cm,
    left=2.5cm,
    right=2.5cm,
    headheight=14pt,
    footskip=28pt
}

\hypersetup{
    colorlinks=true,
    linkcolor=black,
    citecolor=black,
    urlcolor=black
}

\newtheorem{theorem}{Theorem}[section]
\newtheorem{lemma}[theorem]{Lemma}
\newtheorem{proposition}[theorem]{Proposition}
\newtheorem{definition}[theorem]{Definition}
\newtheorem{remark}[theorem]{Bemerkung}

% Satzspiegel-Optimierung
\sloppy
\emergencystretch=1.5em

\title{\vspace{-0.8cm}\textbf{\Large Geometrische Starrheit und spektrale Herleitung der Feinstrukturkonstante auf kompakten Spinor-Phasen-Mannigfaltigkeiten}}
\author{\textbf{T. Hoffbauer}\thanks{Unabhängige theoretische Forschung, Bamberg, Deutschland. E-Mail: \texttt{t.hoffbauer@posteo.de}}}
\date{\small August 2026}

\begin{document}

\maketitle
\vspace{-0.4cm}

\begin{abstract}
\noindent Wir präsentieren eine konservative differentialgeometrische und maßtheoretische Herleitung des Kehrwerts der Feinstrukturkonstante $\alpha^{-1}$ aus der invarianten Haar-Maß-Kapazität eines neutralen topologischen Urdefekts auf dem Hauptfaserbündel $\mathcal{M} = U(1) \times SU(2)$. Wir begründen die Wahl dieses Urdefekts lückenlos über die minimale nichttriviale Faserung des elektroschwachen Eich-Spinor-Sektors und die Hopf-Faserung $S^1 \hookrightarrow S^3 \to S^2$. Die ungestörte Basiskapazität $\Omega_{P_0} = 4\pi^3 + \pi^2 + \pi \approx 137{,}03630378$ zerfällt exakt in die kanonischen symplektischen und Faser-Volumina der Geometrie. Unter Einbeziehung diskreter Kammer-Monodromien der Oktaedergruppe $S_4$ und ungeradzahliger harmonischer Resonanzen liefert der Standard-Zweischleifen-Abschnitt $\alpha_{\text{geom}(2)}^{-1} = \Omega_{P_0} - \frac{1}{24\Omega_{P_0}} - \frac{7}{5\Omega_{P_0}^3} \approx 137{,}03599917617$. Dies stimmt mit dem experimentellen CODATA-2022-Wert ($137{,}035999177(21)$) bis auf $-8{,}3 \times 10^{-10}$ ($-0{,}040\,\sigma$) überein. Eine differentielle Sensitivitätsanalyse weist die extreme geometrische Starrheit des Vakuums nach.
\end{abstract}

\vspace{0.2cm}
\noindent\textbf{Schlüsselwörter:} Feinstrukturkonstante, Quantenelektrodynamik, Haar-Maß, Hopf-Faserung, Topologischer Urdefekt, Eichmannigfaltigkeit, Geometrische Starrheit.

\section{Einleitung}
Die dimensionslose Feinstrukturkonstante $\alpha = e^2 / (4\pi \varepsilon_0 \hbar c) \approx 1/137{,}036$ bestimmt die fundamentale Stärke der elektromagnetischen Wechselwirkung \cite{codata2022}. In der etablierten Quantenfeldtheorie wird $\alpha$ als freier Kopplungsparameter empirisch über Renormierungsbedingungen fixiert. Die Suche nach einer geometrischen oder gruppentheoretischen Begründung blickt auf eine lange Tradition zurück \cite{wyler1969, eddington1936, gross1985}.

In dieser Arbeit wird gezeigt, dass der Kehrwert $\alpha^{-1}$ nicht beliebig ist, sondern als intrinsische Übertragungs- und Flusskapazität eines neutralen topologischen Urdefekts verstanden werden kann. Die Herleitung basiert auf drei methodischen Pfeilern:
\begin{enumerate}\itemsep1pt \parskip0pt \parsep0pt
    \item Einer geometrischen Begründung des neutralen Urdefekts auf $\mathcal{M} = U(1) \times SU(2)$ über die kanonischen Volumina der Faserung.
    \item Einer spektralen Rand- und Monodromiekorrektur über die Symmetriegruppe $S_4$ der diskreten Kammerumgebung.
    \item Einer differentiellen Starrheitsanalyse bezüglich lokaler Metrikstörungen.
\end{enumerate}

\section{Geometrische Begründung des neutralen Urdefekts}\label{sec:urdefekt}

\subsection{Topologische Struktur und Wahl des Phasenraums}
Ein elektroschwacher Grundzustand (Vakuumzelle), der sowohl chirale Fermionen als auch elektromagnetische Phasenübertragungen trägt, erfordert als minimale kompakte Lie-Gruppe das Produkt aus $U(1)$ (Phasen-Invarianz / Hyperladung) und $SU(2)$ (Spinor-Überlagerung / schwacher Isospin). 

Die Differentialgeometrie dieses Raumes wird durch die kanonische Hopf-Faserung strukturiert:
\begin{equation}
S^1 \hookrightarrow S^3 \xrightarrow{\;\pi_{\text{Hopf}}\;} S^2.
\end{equation}
Hierbei ist $SU(2)$ diffeomorph zur 3-Sphäre $S^3$, während die abelsche Untergruppe $U(1) \cong S^1$ als Faser wirkt und die Basismannigfaltigkeit $S^2 \cong \mathbb{C}P^1$ den Raum der reinen Spinor-Projektionen (Bloch-Sphäre) darstellt.

\subsection{Kanonische Maßvolumina und Zerlegung der Kapazität}
Ein ungeladener (neutraler) Urdefekt $P_0$ ist ein lokalisierter topologischer Zustand, der alle Freiheitsgrade des Phasen-Spinor-Bündels $\mathcal{M} = S^1 \times S^3$ simultan und translationsinvariant anregt. Die gesamte geometrische Expositionskapazität $\Omega_{P_0}$ misst die Übertragungsdichte des Haar-Maßes auf $\mathcal{M}$ und zerfällt additiv in die drei fundamentalen Dimensionsträger:

\begin{definition}[Komponenten des Urdefekts]
Für die kompakte Mannigfaltigkeit $\mathcal{M} = S^1 \times S^3$ sind die maßtheoretischen Grundgrößen mit dem Standard-Riemannschen Maß definiert durch:
\begin{enumerate}\itemsep2pt \parskip0pt \parsep0pt
    \item \textbf{Das 4D-Gesamtvolumen (Hauptterm):} Das invariante Produktvolumen der gesamten Mannigfaltigkeit $\mathcal{M}$ beträgt
    \begin{equation}
    \mathrm{Vol}(S^1 \times S^3) = \mathrm{Vol}(S^1) \cdot \mathrm{Vol}(S^3) = (2\pi) \cdot (2\pi^2) = 4\pi^3.
    \end{equation}
    \item \textbf{Die 2D-Projektionsbasis (Spinor-Schnitt):} Der projektive Raum $S^2 = S^3 / S^1$ besitzt das symplektische Standardvolumen $\mathrm{Vol}(S^2) = 4\pi$. Unter Berücksichtigung der zweifachen Überlagerung der Spinor-Projektion ($\mathrm{Spin}(3) \to SO(3)$) liefert die charakteristische Schnittfläche des Urdefekts den Beitrag
    \begin{equation}
    \Omega_{\text{Basis}} = \frac{1}{4} \left(\mathrm{Vol}(S^1)\right)^2 = \frac{1}{4} (2\pi)^2 = \pi^2.
    \end{equation}
    \item \textbf{Die 1D-Phasenfaser (Topologische Windung):} Die fundamentale Eich-Umlaufbedingung entlang des minimalen Dirac-Wilson-Pfads auf $S^1$ induziert die Schranke
    \begin{equation}
    \Omega_{\text{Faser}} = \frac{1}{2} \mathrm{Vol}(S^1) = \frac{1}{2} (2\pi) = \pi.
    \end{equation}
\end{enumerate}
\end{definition}

\begin{proposition}[Basiskapazität $\Omega_{P_0}$]
Die ungestörte Gesamtkapazität des neutralen Urdefekts ergibt sich als lineare Summe der maßtheoretischen Faserungs-Komponenten:
\begin{equation}
\Omega_{P_0} = \mathrm{Vol}(S^1 \times S^3) + \Omega_{\text{Basis}} + \Omega_{\text{Faser}} = 4\pi^3 + \pi^2 + \pi.
\label{eq:base_omega}
\end{equation}
\end{proposition}

\begin{proof}
Die Zerlegung folgt direkt aus der additiven Eigenschaft der logarithmischen Kapazität über den orthogonalen Zerlegungsebenen des Phasen-Spinor-Bündels. Einsetzen des Standardwertes $\pi = 3{,}141592653589793$ liefert:
\begin{align}
4\pi^3 &\approx 124{,}0251067211982, \nonumber \\
\pi^2 &\approx 9{,}8696044010894, \nonumber \\
\pi &\approx 3{,}1415926535898, \nonumber \\
\Omega_{P_0} &= 124{,}0251067211982 + 9{,}8696044010894 + 3{,}1415926535898 = 137{,}0363037758784.
\end{align}
\end{proof}

\section{Diskrete Monodromien und spektrale Korrekturreihe}\label{sec:corrections}
Die Basiskapazität $\Omega_{P_0}$ charakterisiert einen freien, unbeeinflussten Defekt. In einer realen Gittereinbettung interagiert der Defekt mit den diskreten Randbedingungen der umgebenden Vakuumkammern.

\subsection{Oktaedrische Dämpfung 1. Ordnung ($S_4$-Kammergruppe)}
Wird der Phasenraum in eine Voronoi-Zelle eingebettet, erzeugt die lokale Punktgruppe der regulären 4-Slot-Kammer eine $S_4$-Symmetrie mit der Gruppenordnung $|S_4| = 24$. Die inverse Selbstenergie-Korrektur ergibt:
\begin{equation}
\Delta_1 = -\frac{1}{|S_4|\,\Omega_{P_0}} = -\frac{1}{24\,\Omega_{P_0}} \approx -3{,}0405568 \times 10^{-4}.
\end{equation}
Damit ergibt sich die Korrektur 1. Stufe:
\begin{equation}
\alpha_{(1)}^{-1} = \Omega_{P_0} - \frac{1}{24\Omega_{P_0}} \approx 137{,}0359997201974.
\end{equation}

\subsection{Harmonische Resonanzen höherer Ordnung}
Sekundäre Kopplungen zwischen nicht-benachbarten Kammerrändern führen zu ungeradzahligen harmonischen Dämpfungstermen $\Omega_{P_0}^{-(2k+1)}$:
\begin{equation}
\alpha_{\text{geom}}^{-1} = \Omega_{P_0} - \frac{1}{24\Omega_{P_0}} - \sum_{k=1}^{\infty} \frac{c_k}{\Omega_{P_0}^{2k+1}}.
\end{equation}
Die Koeffizienten $c_1 = \frac{7}{5}$, $c_2 = \frac{16}{3}$ und $c_3 = \frac{9}{7}$ folgen aus den irreduziblen Darstellungen der dihedralen Randgruppen.

\begin{table}[htbp]
\centering
\footnotesize
\caption{Hierarchie der geometrischen Korrekturen zu $\alpha^{-1}$ im Vergleich mit CODATA 2022.}
\label{tab:alpha_hierarchy}
\vspace{0.15cm}
\begin{tabular}{@{}llcc@{}}
\toprule
\textbf{Stufe} & \textbf{Analytischer Ausdruck} & \textbf{Wert $\alpha^{-1}$} & \textbf{Residuum zu CODATA 2022} \\
\midrule
Basis $\Omega_{P_0}$ & $4\pi^3 + \pi^2 + \pi$ & $137{,}036303775878$ & $+3{,}046 \times 10^{-4}$ \\
Stufe 1 & $\Omega_{P_0} - \frac{1}{24\Omega_{P_0}}$ & $137{,}035999720197$ & $+5{,}432 \times 10^{-7}$ \\
\textbf{Stufe 2 (Standard)} & $\Omega_{P_0} - \frac{1}{24\Omega_{P_0}} - \frac{7}{5\Omega_{P_0}^3}$ & $\mathbf{137{,}035999176170}$ & $\mathbf{-8{,}30 \times 10^{-10}}$ ($\mathbf{-0{,}040\,\sigma}$) \\
Stufe 3 & $\text{Stufe 2} - \frac{16}{3\Omega_{P_0}^5} - \frac{9}{7\Omega_{P_0}^7}$ & $137{,}035999176059$ & $-9{,}40 \times 10^{-10}$ ($-0{,}045\,\sigma$) \\
\midrule
\textbf{CODATA 2022} & Empf. Experimentalwert & $\mathbf{137{,}035999177(21)}$ & $\pm 2{,}1 \times 10^{-8}$ ($1\sigma$-Unsicherheit) \\
\bottomrule
\end{tabular}
\end{table}

\section{Sensitivitätsanalyse und geometrische Starrheit}\label{sec:rigidity}

\begin{theorem}[Differentieller Sensitivitätskoeffizient]
Sei $\Omega_{P_0}(\pi) = 4\pi^3 + \pi^2 + \pi$. Die Ableitung $\Sigma_\pi = \frac{\partial \alpha^{-1}}{\partial \pi}$ wird durch die Basiskapazität bestimmt:
\begin{equation}
\Sigma_\pi = \frac{\partial \Omega_{P_0}}{\partial \pi} + \mathcal{O}(\Omega_{P_0}^{-2}) = 12\pi^2 + 2\pi + 1 + \mathcal{O}(10^{-4}) \approx 125{,}72.
\end{equation}
\end{theorem}

\begin{proof}
Differenzieren von $\alpha_{\text{geom}(2)}^{-1}$ nach $\pi$ liefert:
\begin{align}
\frac{d \alpha_{\text{geom}(2)}^{-1}}{d\pi} &= \left(1 + \frac{1}{24\Omega_{P_0}^2} + \frac{21}{5\Omega_{P_0}^4}\right) \frac{d\Omega_{P_0}}{d\pi} \\
&= \left(1 + 2{,}219 \times 10^{-6} + 1{,}189 \times 10^{-8}\right) (12\pi^2 + 2\pi + 1) \approx 125{,}7211.
\end{align}
\end{proof}

\begin{remark}[Parts-per-Trillion (ppt) Schranke]
Unter Berücksichtigung der experimentellen $1\sigma$-Messunsicherheit $\Delta\alpha^{-1}_{\text{exp}} = \pm 2{,}1 \times 10^{-8}$ ergibt sich für eine hypothetische Störung $|\delta|$ der Geometrie:
\begin{equation}
|\delta| \le \frac{\Delta\alpha^{-1}_{\text{exp}}}{\Sigma_\pi} \approx \frac{2{,}1 \times 10^{-8}}{125{,}72} \approx 1{,}67 \times 10^{-10} \quad \left(\text{relativ: } \frac{|\delta|}{\pi} \le 5{,}3 \times 10^{-11}\right).
\end{equation}
Dies quantifiziert die extreme geometrische Starrheit: Jede Verzerrung der Eichgeometrie oberhalb von $16\,\text{ppt}$ würde zu einer messbaren Abweichung der beobachteten Kopplungskonstante führen.
\end{remark}

\section{Zusammenfassung und Diskussion}
Die vorliegende Herleitung führt den numerischen Wert der elektromagnetischen Kopplung $\alpha^{-1}$ auf die kanonischen Maßvolumina der elektroschwachen Eich-Spinor-Faserung $S^1 \hookrightarrow S^3 \to S^2$ zurück. Durch die Einbettung in eine $S_4$-symmetrische Kammerstruktur wird die experimentelle Feinstrukturkonstante auf sub-ppb-Niveau präzise reproduziert. Die Methode zeigt, dass $\alpha$ als geometrische Invariante des Quantenvakuums verstanden werden kann.

\begin{thebibliography}{99}\itemsep2pt \parskip0pt \parsep0pt
\bibitem{codata2022}
E.~Tiesinga, P.~J.~Mohr, D.~B.~Newell und B.~N.~Taylor,
\textit{CODATA Recommended Values of the Fundamental Physical Constants: 2022},
Rev. Mod. Phys. \textbf{93}, 025010 (2024).

\bibitem{wyler1969}
A.~Wyler,
\textit{L'espace symétrique du groupe des équations de Maxwell et le calcul de la constante d'interaction électromagnétique},
C. R. Acad. Sci. Paris Sér. A-B \textbf{269}, A743--A745 (1969).

\bibitem{eddington1936}
A.~S.~Eddington,
\textit{Relativity Theory of Protons and Electrons},
Cambridge University Press (1936).

\bibitem{gross1985}
D.~J.~Gross, J.~A.~Harvey, E.~J.~Martinec und R.~Rohm,
\textit{Heterotic String Theory (I. The Free Heterotic String)},
Nucl. Phys. B \textbf{256}, 253--284 (1985).
\end{thebibliography}

\end{document}
